\section*{四、实验探讨与思考题}

\subsection*{1. 实验心得}

通过本次实验，我学会了利用自组惠斯登电桥测量未知电阻，也了解了非平衡电桥将电阻变化转换为电压信号的原理。在实验过程中，应根据待测电阻的大致范围选择合适的比率臂，并尽量调大电桥灵敏度；调平时观察检流计的偏转方向，再逐步逼近平衡点，计算 $R_x$ 时保留4位有效数字；测量 Cu50 时，应等待温度和电压示数都稳定后再同时读数，以减小热滞后的影响。此外，实验记录不能只保留测量数值，还要记下仪器的误差限、实际工作以及初始化条件，便于后续的不确定度和误差分析。

\begingroup
\setstretch{1.30}
\subsection*{2. 思考题}

\begin{enumerate}
  \item \textbf{为了提高电桥测量灵敏度，应采取哪些措施？为什么？}在不过载的前提下，可以适当提高电源电压，选用更灵敏的检流计，并使各桥臂阻值尽量接近。本实验中，减小比率臂的绝对值也能提高灵敏度，使微小电阻变化引起更明显的指针偏转。

  \item \textbf{用电桥测电阻时，若线路接通后检流计指针总是往一个方向偏转或总不偏转，试分析是什么原因？}指针总向一侧偏转，说明调节过程中没有越过平衡点，可能是比率臂不合适、电阻箱量程不够或接线有误。指针始终不动，则应检查电源、检流计回路和各接线点是否正常。

  \item \textbf{惠斯登电桥比率臂选取的原则是什么？为什么要这样选取？}应先估计待测电阻，再选择比率臂，使平衡时的标准电阻落在电阻箱的有效量程内。条件允许时尽量取 $k\approx1$，这样既能充分利用有效数字，也有利于提高灵敏度。

  \item \textbf{简述非平衡电桥与平衡电桥之间的区别。}平衡电桥通过调节桥臂使输出为零，再由平衡条件求未知电阻；非平衡电桥则根据输出电压的大小求电阻变化，更适合连续测量。

  \item \textbf{非平衡电桥在工程中有哪些应用？举例说明。}它能把微小的电阻变化转换为电压信号，因此常与电阻式传感器配合使用。例如，铜电阻电桥可用于测温，应变片电桥可用于电子秤、压力测量和结构应变监测。
\end{enumerate}
\endgroup
